\documentclass[11pt]{article}
\usepackage[T1]{fontenc}
\usepackage[a4paper,margin=1in]{geometry}
\usepackage{amsmath,amssymb,amsthm}
\usepackage[hidelinks]{hyperref}
\usepackage{microtype}

\newtheorem*{question}{ArXiv-version issue}
\newtheorem*{answer}{Published answer}

\setlength{\parindent}{0pt}
\setlength{\parskip}{0.55em}

\title{Literature Status: The Extra Projection Condition\\
in arXiv:2010.05726 Is Unnecessary}
\author{Run \texttt{fa\_banach\_001} --- lane 1}
\date{13 August 2026}

\begin{document}
\maketitle

\begin{center}
\fbox{\parbox{0.9\textwidth}{
\textbf{Status: answered in the published revision.}
The arXiv version imposes an extra condition for strong convergence of
the metric projections and says its necessity needs investigation.  The
2022 journal version removes that condition completely and proves the
unconditional result.
}}
\end{center}

\section{Issue in the arXiv version}

In arXiv:2010.05726v2, A. B\"erd\"ellima proves that the iterates
$x_n=T x_{n-1}$ of a nonexpansive, quasi-$\alpha$-firmly nonexpansive
map on a complete $\operatorname{CAT}(0)$ space converge weakly to some
$x^*\in\operatorname{Fix}T$~\cite{arxiv}.  Theorem~3 (arXiv PDF pages
11--12) asserts strong convergence
\[
 P_{\operatorname{Fix}T}x_n\longrightarrow x^*
\]
only under condition (34).  The cyclic-projection Theorem~4 on arXiv
PDF page~13 similarly requires condition (40).  The historical remarks
on that page state that the necessity of (40) needs to be investigated.

\begin{question}
Is the additional condition actually needed to conclude that the
metric projections converge strongly to the weak limit?
\end{question}

\section{Published answer}

The revised paper appeared as \emph{On a notion of averaged mappings in
$\operatorname{CAT}(0)$ spaces}, \emph{Functional Analysis and Its
Applications} \textbf{56} (2022), 27--36~\cite{published}.  Its
Theorem~3 (Russian original, journal pages 46--47; PDF pages 9--10)
states the same general convergence theorem with no extra condition:

\begin{answer}
If $T$ is nonexpansive and quasi-$\alpha$-firmly nonexpansive on a
complete $\operatorname{CAT}(0)$ space, then $x_n=T x_{n-1}$ converges
weakly to some $x^*\in\operatorname{Fix}T$, and
$P_{\operatorname{Fix}T}x_n$ converges strongly to $x^*$.
\end{answer}

Theorem~4 (journal pages 47--48) gives the corresponding unconditional
statement for cyclic projections, and Theorem~5 (pages 48--49) does so
for averaged projections.

\section{Why the condition disappears}

The published proof supplies a short completion of the arXiv argument.
The following formulation isolates the decisive step.  Put
$C=\operatorname{Fix}T$ and $p_n=P_Cx_n$.  Since $(x_n)$ is Fej\'er
monotone with respect to $C$, for $m\ge n$ the projection inequality
gives
\[
 d(p_m,p_n)^2
 \le d(x_n,p_n)^2-d(x_m,p_m)^2.
\]
The nonincreasing sequence $d(x_n,C)$ has a limit, so $(p_n)$ is Cauchy
and $p_n\to p\in C$.

Because $x^*\in C$, the same projection inequality yields
\[
 d(x_n,p_n)^2+d(p_n,x^*)^2\le d(x_n,x^*)^2.
\]
Also
$|d(x_n,p)-d(x_n,p_n)|\le d(p,p_n)\to0$.  Hence
\[
 \liminf_n d(x_n,p)\le \liminf_n d(x_n,x^*).
\]
But bounded weak convergence in a complete $\operatorname{CAT}(0)$
space satisfies the Opial inequality: if $p\ne x^*$, then the reverse
strict inequality holds.  Thus $p=x^*$, proving
$P_Cx_n\to x^*$ without any additional regularity condition.

This also shows directly why the conditional arXiv proof was only one
step short: its Cauchy argument already produced $p$; the Opial property
identifies that limit with $x^*$.

\section{Provenance}

A direct attack rediscovered both a geodesic estimate and the Opial
completion before the published-version check.  The exact theorem is,
however, already present in the author's revised publication, so no new
mathematical result is claimed here.  The broad sentence in the arXiv
paper leaving a possible metric notion of monotone operators for future
research is not a precise theorem-sized target and is not promoted.

\begin{thebibliography}{9}
\bibitem{arxiv}
A. B\"erd\"ellima,
\emph{On a notion of averaged operators in $\operatorname{CAT}(0)$ spaces},
arXiv:2010.05726v2 (2020).

\bibitem{published}
A. B\"erd\"ellima,
\emph{On a notion of averaged mappings in $\operatorname{CAT}(0)$ spaces},
Funct. Anal. Appl. \textbf{56} (2022), 27--36;
Russian original: Funktsional. Anal. i Prilozhen. \textbf{56}:1 (2022),
37--50; DOI: 10.1134/S0016266322010038 and 10.4213/faa3859.
\end{thebibliography}

\end{document}
